\section{The Claim Under Test}

There is a thing many Christians believe they know about Easter morning, and
it is usually put in one sentence: \emph{St.\ Peter knew that Jesus had risen
because he noticed how the burial cloths had been left.} The cloth that had
been over the face was not lying with the rest. It was somewhere else in the
tomb, by itself, and something about the way it lay told Peter what had
happened.

The sentence is attractive because it is the kind of detail no one invents. It
has the texture of a thing seen. And its scriptural warrant looks immediate:
John says, of the tomb Peter entered, that there lay ``the linen cloths'' and
also ``the napkin that had been about his head, not lying with the linen
cloths, but apart, wrapped up into one place'' (Jn 20:6--7). Someone noticed
where the face-cloth was. Someone thought it worth reporting sixty years
later. Whoever wrote that verse was reporting a fact about the disposition of
objects in a rock tomb, and reporting it because it mattered.

What the sentence adds to the verse is a name and a conclusion: that the
someone was Peter, and that what he drew from what he saw was the
Resurrection. Both halves are testable, and the Gospels test them
differently. This study takes the sentence as its object. Not the Resurrection
---which is not in question here---but this particular claim about who
inferred it, and from what, and whether the arrangement of the cloths carries
the evidentiary weight that two very different traditions of argument have
placed on it.

Three things have to be worked, and they are not equally difficult.

The first is the text itself. John 20:3--10 is eight verses of Greek in which
almost every word that matters is doing something a translation flattens.
\emph{Othonia} and \emph{soudarion} are not ``linen cloths'' and ``napkin'' in
any sense a modern reader's furniture supplies. \emph{Keimena} says the cloths
were lying; the participle applied to the face-cloth is not \emph{keimenon} but
\emph{entetyligmenon}, and the difference between those two words is the whole
of the observation. \emph{Ch\=oris} and \emph{eis hena topon} locate it. The
Latin the Church read for fifteen centuries makes some of these decisions
differently, and the English most Catholic readers know was made from that
Latin. Section~2 sets out the passage in the three languages; section~3 works
the vocabulary.

The second is the sentence's own difficulty, and it is the centre of this
study. John does not say that Peter believed. John says that the other
disciple believed. Peter goes in first, and sees the most, and the evangelist
gives him the fullest description of the arrangement in the whole
passage---and then attributes the believing, in the next verse, to somebody
else. Luke, telling what is evidently a version of the same episode, has Peter
look in, see the linen cloths lying by themselves, and go away
\emph{wondering}. Neither evangelist says what the popular sentence says. The
question is what they do say, whether the harmonizations offered since
antiquity survive inspection, and what remains of the claim afterwards.
Sections~4 and~5 take that up.

The third is the question of why the arrangement should matter at all, and
here two quite separate explanations are in circulation, of very unequal
standing.

One is ancient and genuinely attested. It is an argument against theft: that
grave-robbers do not undress a corpse, and that a body glued into its
wrappings by a hundred pounds of myrrh and aloes is not unwrapped in a hurry by
anyone whose object is the body. John Chrysostom makes this argument at length
in his eighty-fifth homily on John, and it is not a devotional flourish but a
piece of forensic reasoning about what thieves do. Section~6 reads it at its
own locus and asks what it establishes.

The other is modern, and it is the reason many readers have met the sentence at
all. It holds that a folded napkin was a signal in Jewish domestic custom: that
a master rising from table who crumpled his napkin meant that he was finished,
while a master who folded it meant \emph{I am coming back}, and that ``every
Jewish boy'' would have known this. On that reading the face-cloth in the tomb
is not evidence but a message. Section~8 asks whether any such custom is
attested anywhere, and where the claim came from. The finding is reported
there in full; it is not what the claim's circulators suppose.

A word about method, since this study reaches a negative conclusion about a
popular reading and a partly negative conclusion about the maintainer's own
framing of the question. Both were reached by reading, not by preference. Every
witness named here was read at its own locus: Chrysostom in the printed
Nicene and Post-Nicene Fathers volume, not in a quotation; Augustine in the
same series and in the Latin; the Greek in a registered edition of the
Byzantine textform; the Latin in the page images of a Clementine Vulgate
printing. Where a search failed to find something, the search is described and
its boundaries are stated, because a failed search is a finding only to the
extent that its shape is known. Where the tradition disagrees with itself, the
disagreement is left standing. Nothing here is harmonized for tidiness.

Two large subjects are deliberately outside the argument and are set out as a
boundary in section~10: the Shroud of Turin and the Sudarium of Oviedo. They
are a different question with a different literature, and half-treating them
would make this study worse rather than longer.
